% AI-produced by GPT-6 Astra (OpenAI, ChatGPT Pro mode).
% Reconstructed on 2026-11-20 from the complete accepted mathematical argument
% and in-session root/independent audits. This is not a claim of recovering
% the lost historical TeX bytes. See logs/A144_A146_recovery_provenance.md.
\Needspace{10\baselineskip}
\subsection{A.144. Realizing the obstruction for finite fundamental groups}

\begin{saresult}{A.144}{Proved under stated hypotheses}{Surjectivity of the three-dimensional component obstruction for finite fundamental group}
Let $B$ be a closed, connected, oriented real PL three-manifold with
finite fundamental group. The homomorphism of A.143 is surjective:
\[
 c_B:\pi_0P_{\mathrm{TOP}}(B)\longrightarrow
                  H^2(B;\mathbb Z/2).
 \tag{144.1}
\]
Every class on the right is represented by a TOP concordance that is
the identity on both ends. The image of the PL component map, and the
image of the definable component map for every fixed real o-minimal
expansion, lie in its kernel. No equality of either image with this
kernel is asserted.
\end{saresult}

\begin{proof}
We first give the relative realization needed both here and in A.145.
Let $X$ be a compact, connected, oriented smooth four-manifold with
finite fundamental group $\Gamma$. Assume that a Sylow two-subgroup
$\rho\subset\Gamma$ satisfies
\[
                 SK_1(\mathbb Z[\rho])=0.
 \tag{144.2}
\]
For every $\eta\in H^3(X,\partial X;\mathbb Z/2)$ there is a
homeomorphism $f:X\to X$, equal to the identity on the entire boundary,
with formal relative PL obstruction $c_{\mathrm{PL}}(f)=\eta$.

Put $Z=X\times I$, with its full boundary. The relative normal-invariant
coordinates and universal PL-obstruction map used in A.142 are
\[
 \begin{aligned}
 [Z/\partial Z,G/\mathrm{TOP}]
   &\cong H^2(Z,\partial Z;\mathbb Z/2)
                  \oplus H^4(Z,\partial Z;\mathbb Z),\\
 \xi_*(x,y)&=\operatorname{red}_2(y)
             \quad\text{on this suspension pair}.
 \end{aligned}
 \tag{144.3}
\]
In particular the formula is valid at $x=0$, which is all we use.
The source and the boundary convention are
\cite[Lemmas 4.8--4.9, printed p.~2526]{Galvin2026};
the full relative interpretation, including the side labeling, was
proved in (G.9)--(G.14). In general the formal formula can have a
$\operatorname{Sq}^2x$ term. Here that term vanishes because $x$
is the suspension of a degree-one relative class, and Steenrod
squares commute with suspension.

We do not claim that $L^s_5(\mathbb Z[\Gamma])$ vanishes. Instead,
apply the obstruction formula of
\cite[Theorem A(b), printed p.~351]{HambletonMilgramTaylorWilliams1988}.
Its hypotheses concern finite Poincare pairs and the decoration
$U\subset\operatorname{Wh}(\Gamma)$. It requires that
\[
 Y=\operatorname{im}\bigl(SK_1(\mathbb Z[\rho])
                   \longrightarrow\operatorname{Wh}(\Gamma)\bigr)
                         \subset U.
\]
By (144.2) we may choose $U=0$, so this is the \emph{simple} surgery
obstruction formula. The reference normal map is the identity of $Z$.
In dimension five the formula has the form
\[
       \theta(\nu)=\kappa_3^s\bigl(c_*
                   \operatorname{ARF}_3(\nu)\bigr).
 \tag{144.4}
\]
Here $c:Z\to B\Gamma$ classifies the universal cover. The degree-three
homology component in the source's definition (0.6) can only come
from its degree-two universal class and the degree-zero Wu factor.
Thus, in the coordinates (144.3),
\[
       \operatorname{ARF}_3(x,y)
              =x\mathbin{\cap}[Z,\partial Z],
       \qquad \theta(0,y)=0.
 \tag{144.5}
\]
The integral degree-four coordinate does not contribute to
$\operatorname{ARF}_3$. This uses the actual relative-pair theorem,
not only a closed-manifold corollary.

The integral coefficient reduction in degree four is onto:
\[
 H^4(Z,\partial Z;\mathbb Z)\longrightarrow
 H^4(Z,\partial Z;\mathbb Z/2)\longrightarrow
 H^5(Z,\partial Z;\mathbb Z)
            \xrightarrow{\,2\,}H^5(Z,\partial Z;\mathbb Z).
 \tag{144.6}
\]
The last group is $\mathbb Z$ by relative orientation. Choose $y$
whose reduction is the interval suspension $\sigma\eta$.
Equation (144.5) makes the simple surgery obstruction of $(0,y)$
zero. Relative TOP surgery realizes this normal invariant by a
simple homotopy equivalence
\[
 \phi:(N;\partial_+N,\partial_0N,\partial_-N)
 \longrightarrow
 (X\times I;X\times0,\partial X\times I,X\times1)
 \tag{144.7}
\]
whose full boundary map is the prescribed homeomorphism.
Finite groups are good for the relative four-dimensional TOP
$s$-cobordism theorem; this is the same good-group input used in the
relative product argument of A.142. The resulting five-dimensional
$s$-cobordism has a product homeomorphism
$p:N\to\partial_+N\times I$ which fixes the bottom and the labeled
side product. This is precisely the relative conclusion in
\cite[Proposition 4.10 and its diagram, printed pp.~2526--2527]
{Galvin2026}. With the same notation as (G.11)--(G.13), the map
\[
                  f=\phi_+p_-\phi_-^{-1}
 \tag{144.8}
\]
is the identity on $\partial X$: both boundary labelings agree there
because $p$ fixes the whole side product. The induced
mapping-cylinder identification preserves the prescribed boundary
PL reductions. Consequently (144.3) gives
$c_{\mathrm{PL}}(f)=\eta$. This proves the relative realization
assertion.

We now verify (144.2) for the groups in the theorem without invoking
geometrization or a classification of finite three-manifold groups.
Choose a finite CW triangulation of $B$. Since $\Gamma$ is finite,
the universal cover $\widetilde B$ is a compact, connected, simply
connected, oriented three-manifold. Its homology is
\[
 H_0(\widetilde B;\mathbb Z)=H_3(\widetilde B;\mathbb Z)=\mathbb Z,
       \qquad H_1(\widetilde B;\mathbb Z)
              =H_2(\widetilde B;\mathbb Z)=0.
 \tag{144.9}
\]
Here $H_1=0$ follows from simple connectivity, and $H_2=0$ follows
from Poincare duality. Every deck transformation preserves the
orientation induced from $B$, so its action on $H_3$ is trivial.

The lifted cellular chain modules $C_i$ are finite free
$\mathbb Z[\Gamma]$-modules, and (144.9) gives an exact sequence of
$\Gamma$-modules
\[
 0\longrightarrow\mathbb Z\longrightarrow C_3
   \longrightarrow C_2\longrightarrow C_1
   \longrightarrow C_0\longrightarrow\mathbb Z
   \longrightarrow0,
 \tag{144.10}
\]
with trivial modules at both ends. Splice copies of this extension,
using the augmentation $C_0\to\mathbb Z$ followed by the injection
$\mathbb Z\to C_3$. Exactness is retained at both joins. This is a
four-periodic free resolution of the trivial module, and proves
periodic group cohomology directly.

The final introductory paragraph on printed p.~352 of
\cite{HambletonMilgramTaylorWilliams1988} explicitly states that the
Sylow two-subgroups of groups with periodic cohomology have trivial
$SK_1$, and that the displayed formulas therefore hold in simple
$L$-groups. It does \emph{not} assert that $SK_1(\mathbb Z[\Gamma])$
itself is zero; the same paragraph notes that this larger group can
have two-torsion. Applying the stated Sylow assertion proves (144.2).

Choose one compatible smoothing of $B$ and the product smoothing of
$X=B\times I$. Apply the relative realization to the suspended class
corresponding to any element of $H^2(B;\mathbb Z/2)$. The resulting
$f$ fixes both ends and represents that value under $c_B$ by A.143.
This proves (144.1). Kernel containment and the natural real
definable conclusion are exactly the invariance and comparison
statements in A.143. Nothing in this proof provides an isotopy of
a zero-obstruction four-dimensional map to a PL map.
\end{proof}

